\section{Užduoties tikslas}
\label{sec:uzduoties_tikslas}

Užduoties tikslas -- sudaryti vieno pasirinkto daugiafaktorinio autentifikavimo (\textit{angl. multi-factor authentication}, MFA) metodo implementavimo planą~\cite{VMA}. Pasirinktas metodas -- \textbf{OAuth 2.0} protokolas kartu su jo tapatybės sluoksniu \textbf{OpenID Connect} (OIDC). Antrasis faktorius -- išorinio tapatybės teikėjo (\textit{angl. identity provider}, IdP) \textit{Google} paskyra.

Plane nagrinėjama interneto aplikacija (serveris -- Node.js su TypeScript), kurioje:
\begin{itemize}
    \item \textbf{pirmasis faktorius} -- aplikacijos slaptažodis (tai, ką vartotojas žino);
    \item \textbf{antrasis faktorius} -- prisijungimas prie iš anksto susietos \textit{Google} paskyros per OAuth 2.0 ir OpenID Connect (tai, ką vartotojas turi).
\end{itemize}

Užduoties reikalavimai ir juos atitinkančios darbo dalys pateiktos \ref{tab:reikalavimai} lentelėje.

\begin{table}[H]
    \centering
    \caption{Užduoties reikalavimai ir juos atitinkančios darbo dalys.}
    \label{tab:reikalavimai}
    \begin{tabularx}{\textwidth}{L l}
        \toprule
        \textbf{Reikalavimas} & \textbf{Darbo dalis} \\
        \midrule
        Pasirinktas MFA metodas ir trumpas aprašymas, kaip jis veikia & \ref{sec:metodas} skyrius, \ref{fig:oidc_srautas} pav. \\
        Kaip vartotojas aktyvuos MFA metodą & \ref{sec:aktyvavimas} skyrius, \ref{fig:aktyvavimas} pav. \\
        Kaip vartotojas autentifikuosis MFA metodu & \ref{sec:autentifikavimas} skyrius, \ref{fig:prisijungimas} pav. \\
        Kokie duomenys saugomi duomenų bazėje & \ref{sec:duomenu_baze} skyrius, \ref{fig:db_schema} pav. \\
        Biblioteka arba servisas, kuriuo implementuojamas MFA & \ref{sec:biblioteka} skyrius \\
        Aktualūs paveikslėliai & \ref{fig:oidc_srautas}--\ref{fig:db_schema} pav. \\
        \bottomrule
    \end{tabularx}
\end{table}

Plano tekstui ir diagramoms parengti buvo pasinaudota Claude~\cite{claude}.
